\begin{song}{Die Gedanken Sind Frei}{Deutsches Volkslied}

    \begin{verse}  
        Die Ge[G]danken sind frei, wer [D7]kann sie er[G]raten? \\
        Sie fliehen vorbei wie [D7]nächtliche [G]Schatten. \\
        Kein [D]Mensch kann sie [G]wissen, kein [D]Jäger er[G]schießen
        mit [C]Pulver und [G]Blei. \\
        Die Ge[D]danken sind [G]frei!
    \end{verse}  

    \begin{verse}  
        Ich [G]denke, was ich will und [D7]was mich be[G]glücket,\\
        doch alles in der Still’
        und [D7]wie es sich [G]schicket.\\
        Mein [D]Wunsch und Be[G]gehren
        kann [D]niemand ver[G]wehren,
        es [C]bleibet da[G]bei:\\
        Die Ge[D]danken sind [G]frei!\\
    \end{verse}  

    \begin{verse}  
        Und [G]sperrt man mich ein
        im [D7]finsteren [G]Kerker, \\
        das alles sind rein
        ver[D7]gebliche [G]Werke, \\
        denn [D7]meine Ge[G]danken
        zer[D]reißen die [G]Schranken
        und [C]Mauern ent[G]zwei. \\
        Die Ge[D]danken sind [G]frei!
    \end{verse}  

    \begin{verse}  
        Drum [G]will ich auf immer
        den [D7]Sorgen ent[G]sagen \\
        und will mich auch nimmer
        mit [D7]Grillen mehr [G]plagen. \\
        Man [D]kann ja im [G]Herzen
        stets [D]lachen und [G]scherzen
        und [C]denken da[G]bei: \\
        Die Ge[D]danken sind [G]frei!
    \end{verse}  

\end{song}
